\section{Scope and Assumptions}
	
	The scope of this threat model is defined around the current and future architecture of \textbf{ViperVault}.
	At present, the application is deployed as a local-only installation on Android and iOS devices.
	All credential data is stored directly on the device and protected through strong cryptographic mechanisms.
	There is no network communication and no persistent cloud storage, which significantly reduces the remote attack surface.
	Looking ahead, however, the project’s roadmap includes several major extensions: the release of desktop clients for Windows, macOS and Linux; the development of browser extensions to facilitate password autofill;  the addition of a cloud synchronization service to enable seamless access across devices.
	These planned features will necessarily expand the system boundaries and introduce new classes of threats that must be considered from the outset.
	In terms of adversaries, the model assumes a broad spectrum of potential threat actors.
	These include opportunistic attackers exploiting common weaknesses, more sophisticated advanced persistent threats (APTs) conducting targeted campaigns, malicious insiders with privileged access and supply-chain adversaries attempting to compromise the build or distribution process.
	The security goals of ViperVault are fourfold.
	First, the confidentiality of all stored credentials must be preserved under every circumstance.
	Second, the integrity of both the vault and its metadata must be ensured, preventing unauthorized tampering.
	Third, the application must remain reliably available to legitimate users.
	Finally, non-repudiation must be enforced for critical operations such as master-password changes, ensuring accountability for sensitive security events.
	The cryptographic foundations of the system rest on well-established and modern primitives.
	The vault is encrypted using the XChaCha20-Poly1305 algorithm, which offers strong authenticated encryption guarantees.
	The master password is processed through Argon2id, a memory-hard key derivation function designed to resist GPU and ASIC cracking attempts.
	All keys are derived exclusively on the user’s device and are never transmitted, in line with zero-knowledge design principles.